\documentclass[11pt]{article}
\usepackage[margin=1.03in]{geometry}
\usepackage{amsmath,amssymb,amsthm}
\usepackage{setspace}\usepackage{textcomp}
\usepackage{booktabs}
\usepackage{graphicx}
\graphicspath{{./figures/}{./}}
\usepackage{array,multirow}
\usepackage[colorlinks=true,linkcolor=blue,citecolor=blue]{hyperref}
\usepackage[round]{natbib}
\setstretch{1.12}
\newcommand{\bp}{\,\mathrm{bp}}
\newcommand{\sBE}{\sigma_{BE}}
\newcommand{\sIV}{\sigma_{IV}}

\title{Two Prices for One Volatility?\\ Convexity across the Curve and the Swaption Market\\[4pt]
\large Paper 2 master --- theory, methods, complete results (v4.5)}
\author{Alessio Ottaviani\thanks{EDHEC Business School. Supervisor: Riccardo Rebonato. This document
consolidates the Project Report (July 2026) and all empirical work through 4 August 2026. Every number
is generated by saved, re-runnable scripts referenced in the Register (Section~\ref{sec:register});
none is quoted from memory. Bibliography verified verbatim against the papers on file.}}
\date{August 4, 2026}

\begin{document}
\maketitle

\begin{abstract}
\noindent The yield curve prices its own convexity: a value- and duration-matched butterfly earns a
negative carry-and-roll whose magnitude reveals a curve-implied break-even volatility, $\sBE$. This
paper extends \citet{rr2021} beyond its four swap markets --- adding four government curves on
independent Bloomberg data --- and answers the four questions its authors
pose in their closing section and leave ``for future analysis'' for lack of reliable swaption prices:
the cost of reproducing the long-gamma profile with a swaption straddle; whether the implicit premium
from negative carry exceeds the option premium; which venue prices convexity more cheaply at a given
moment; and whether a delta-, gamma- and vega-neutral package of swaps and straddles carries a
non-zero cost. \textbf{(i) Segmentation.} The correlation of monthly \emph{changes} between $\sBE$ and
swaption ATM normal vols is $-0.01$ (USD swap), $+0.03$ (UST), $-0.07$ (GBP), $+0.24$ (EUR), $+0.31$
(JPY); across the full swaption cube --- $16$ cells per market, $80$ in total --- $56$ cells ($70\%$)
have $|\rho|<0.20$, and the ten-year tail column is not systematically higher than the others, so the
result is not an artefact of cell selection. A rolling 60-month window shows no closing trend
(USD $-0.35$, GBP $-0.43$, UST $-0.21$ per decade: if anything, \emph{more} segmented).
\textbf{(ii) Within-currency identification.} Adding the Bund, gilt and JGB curves gives eight markets
and, for each currency, two curves priced against one swaption surface. The negative co-movement in
sterling is three times stronger in swaps ($-0.61$) than in gilts ($-0.18$) --- where LDI pension
hedging is executed \citep{ks2023}; in Japan the most integrated curve is the JGB ($+0.79$ vs.\ $+0.65$)
--- the market where a single official actor has been the dominant marginal participant. Segmentation tracks \emph{where the constrained
clientele physically operates}, with institutional differences held constant by construction.
\textbf{(iii) Mechanism.} Timed convexity returns are monotone in dealer-CDS stress terciles in four
of four swap markets (HIGH cell $t=2.0$--$2.6$), and remain significant excluding all four crisis
windows in three of four. Sorting instead on the \citet{hkm2017} intermediary capital ratio --- an
academic measure independent of our CDS composite --- the premium concentrates where intermediaries are
constrained in four of four markets ($t=1.7$--$3.4$). \textbf{(iv) Limits of arbitrage, priced.} With
dealer quotes (full spread $0.10$--$0.15/0.30/1.00\bp$ by regime), net returns are $+7.38$ (USD),
$+10.41$ (EUR), $+8.88$ (GBP), $+4.82$ (JPY) $\bp$/month and the pessimistic full-replacement bound is
positive everywhere: \emph{costs are not the barrier}, which strengthens the balance-sheet mechanism by
exclusion. \textbf{(v) Not a known factor.} Alpha survives the tradable term factors, the DV01-neutral
futures butterfly --- the closest tradable analogue of the trade --- and the \citet{fh2001}
trend-following straddles, with $R^2$ of $0.00$--$0.03$ throughout. \textbf{(vi) Balance sheet, not
macro.} On the first principal component of the wedges across seven markets, intermediary proxies give
$\bar R^2=0.584$ and a macroeconomic placebo $0.130$, a contrast of $4.5\times$.
\textbf{(vii) The structural model is rejected, and the rejection is the result.} A transplant of the
certainty-equivalent/barrier framework, over-identified across markets with one free risk-aversion
parameter, under-predicts the premium everywhere and achieves a negative cross-sectional $R^2$; the
premium tracks segmentation (rank $+0.80$) twice as strongly as it tracks the mark-to-market risk the
arbitrageur bears ($+0.40$). The convexity premium is not compensation for arbitrage risk.
\end{abstract}

\paragraph{What this document is.} The definitive master for Paper~2 of the dissertation. It
consolidates in one place: the theoretical framework and the structural model with its verdict; the
complete data inventory with provenance; the methods; every result to date; the
Established-versus-Open ledger; the run Register for all seventeen scripts; and the road to
completion with a target-journal assessment. Its companion is the package \texttt{convexity/}
(scripts \texttt{check\_inputs}, \texttt{01}--\texttt{17}) in the project repository.

\paragraph{Version history.} \emph{v2.0}: restructured story-first; facts re-derived as prediction
battery S1--S5. \emph{v2.1}: robustness and hygiene audit. \emph{v2.2}: made self-contained.
\emph{v2.3}: every attribution to \citet{rr2021} verified against the published paper (JEF 63,
392--413). \emph{v2.4}: dealer-CDS layer rebuilt with an identity-checked dedup rule. \emph{v2.5}:
data consolidated to one Bloomberg hub read directly; dealer composite cleaned; costs implemented and
gated. \emph{v3.0 (4 August 2026)}: \textbf{costs switched on} with dealer quotes and the narrative
revised accordingly; \textbf{the panel extended from five to eight markets} with the government curves,
delivering within-currency identification; the \textbf{full swaption-cube scan} (80 cells); the
\textbf{term-premium vulnerability} discovered, diagnosed and answered; the \textbf{control battery}
against the first paper's tradable factors; the \textbf{third venue} (bond-future implied volatility);
the \textbf{P5 panel} on the design of the first paper's Table~XIX; and the \textbf{structural model},
built, calibrated and \emph{rejected}, with the diagnosis of why. \emph{v4.5 (11 August 2026)}: added
the \textbf{cross-venue strategy} (Section~\ref{sec:strategy}) --- the trade built four ways, including
the antecedent paper's signal-weighted per-trade design with convergence exit, delta-hedged in a
genuine daily Bachelier simulation, costed with real dealer quotes on the rate legs and a conservative
literature bound on the option leg --- delivering the no-trade band ($5$--$7$ bp/month wedge against
$9$--$14$ round-trip) as the quantitative form of the limits-to-arbitrage thesis, and discharging the
vega-neutral-assembly item; everything prior to this section is unchanged from v4.4.

\section{Introduction and the Delimited Claim}\label{sec:intro}

The claim of this paper is deliberately narrow, and it is worth stating what is \emph{not} claimed
before stating what is. We do not claim that the curve systematically under-prices volatility: across
all eight markets the mean wedge between trailing realised variance and curve-implied variance is
\emph{negative}, which is the familiar variance risk premium and not a finding. We do not claim to have
found a costless arbitrage: the trade requires balance sheet, and its returns are volatile
($\sigma \approx 47$--$76\bp$/month against a mean of $5$--$10$). We do not claim that the premium is
predictable in real time in a way that would survive a fund's investment committee: the conditional
ladder that identifies the mechanism uses in-sample thresholds, and its real-time analogue
(Section~\ref{sec:conditional}) is suggestive rather than significant.

What we claim is this. Two venues price the same object --- the volatility of the long rate --- and
their prices move independently. The independence is not a measurement artefact: it holds across the
whole swaption cube, it holds against a third venue (options on bond futures), it does not close over
time, and it is stronger precisely where the constrained clientele operates. The wedge between the
venues carries a premium; that premium is concentrated in states where intermediary balance sheets are
impaired, by three independent measures of impairment; it is not explained by term, curvature, slope,
level, trend-following, or liquidity factors; and it is not compensation for the arbitrageur's own
mark-to-market risk. The last of these is established by building the structural model that would
imply it, and finding that the data reject it.

\section{Theoretical Framework}\label{sec:theory}

\subsection{The object}
A value- and duration-matched butterfly on the swap or government curve is, to first order, a pure
convexity position: its delta and its cash outlay both vanish, so what remains is gamma. Following
\citet{rp2018} and \citet{rr2021}, holding such a fly earns a carry-and-roll $\theta$ that is negative
when the curve is concave in the relevant region; the position is compensated, if at all, by the
realised gamma P\&L. Setting the two equal defines a break-even volatility,
\begin{equation}
\sBE^2 \;=\; \frac{-2\,\theta}{C\,\Delta t},
\label{eq:sbe}
\end{equation}
where $C$ is the net convexity of the fly in the chosen weighting. Equation~\eqref{eq:sbe} is the
curve's own price of volatility: the level of realised volatility at which the fly breaks even. It is
observable from the curve alone, with no option data. As \citet{rr2021} put it, in this construction
``the role of the implied volatility is played by the curvature of the swap curve''.

\subsection{Venues and segmentation}
The same volatility is quoted directly in the swaption market. If capital moved freely between the two
venues, $\sBE$ and the swaption implied volatility $\sIV$ would be tied by arbitrage --- not equal,
since the instruments differ in maturity, moneyness and convexity profile, but co-moving. The central
empirical object of this paper is the joint behaviour of the two. The story we propose for their
independence is \emph{one object, two clienteles}: the curve is shaped by duration-motivated flow
(pension and insurance liability hedging, central-bank operations, index investors) that is
price-insensitive to volatility and does not arbitrage it, while the swaption market is intermediated
by option desks whose inventory and risk limits determine $\sIV$. The two clienteles do not meet, and
the intermediaries who could connect them face a balance-sheet cost.

\subsection{Crisis asymmetry}
The story implies a signature in crises, and it is not one-directional. When volatility spikes and
capital flies to the long end, the curve becomes convexity-\emph{rich} (2008, 2011). When balance
sheets are the binding constraint and duration hedgers are forced sellers, the curve becomes
convexity-\emph{cheap} (2022--23, in the markets where the constraint bound). Both are stress; they
move the wedge in opposite directions. This two-channel structure is a prediction of the story, and it
recurs throughout the results --- most sharply in the P5 panel (Section~\ref{sec:p5}), where funding
and dealer-credit proxies enter with one sign and volatility/illiquidity proxies with the other.

\subsection{The prediction battery}\label{sec:battery}
\begin{description}
\item[S1] The wedge is driven by intermediary balance-sheet conditions, not by measurement noise.
\item[S2] Markets differ in integration, and the ordering follows the institutional structure of the
clientele: Japan most integrated (single dominant actor operating in both), the United Kingdom least
(LDI hedging concentrated in swaps).
\item[S3] Crisis asymmetry: rich in vol-stress episodes, cheap in balance-sheet-stress episodes.
\item[S4] The premium is concentrated in constrained states, not earned uniformly.
\item[S5] Reversal half-lives are longer where segmentation is greater.
\end{description}

\subsection{The structural model, and why it matters that it fails}\label{sec:structural-theory}
The empirical battery establishes that the premium exists and where it lives. It does not explain why
the wedge settles at the level it does, nor why that level differs across markets. To answer that we
transplant the certainty-equivalent/barrier framework that \citet{rebonatoNote} develops for
convergence trades. The arbitrageur who would close the wedge faces three things, not one: slow
convergence, so the expected gain accrues over months; adverse mark-to-market along the path, arriving
precisely when the wedge widens and other arbitrageurs are also losing; and a loss barrier (VaR or
margin) that can force unwinding before convergence, crystallising the loss --- the
\citet{sv1997} channel. He enters only if the certainty equivalent of the P\&L is non-negative. The
equilibrium premium $\mu^\ast$ is the one that makes him indifferent: $CE(\mu^\ast)=0$. Below it nobody
enters; above it entry continues until the premium is competed back.

Two design choices are deliberate. First, utility is \textbf{CARA on the P\&L} rather than CRRA on
wealth: under CRRA the parameter is $\gamma=\mathrm{RRA}/W_0$, so the normalisation by the
arbitrageur's wealth is a stated choice rather than an estimated quantity, whereas absolute risk
aversion under CARA is measured in units of observable P\&L and requires no wealth to be postulated.
Second, and more important, the model here is \textbf{over-identified}. A framework calibrated on a
single market has one $\lambda^\ast$ and one free risk-aversion parameter, so the parameter absorbs
whatever level is observed and the model cannot fail. Here the dynamics
$(\sigma,\rho)$ are estimated separately in each market and one risk-aversion parameter must reproduce
the whole cross-section of premia. That is a testable restriction, and Section~\ref{sec:structural}
reports that the data reject it.

\section{Related Literature}\label{sec:lit}
The construction of $\sBE$ is due to \citet{rp2018} and is applied to four swap markets by
\citet{rr2021}, whose conclusion explicitly leaves the comparison with option-implied volatility ``for
future analysis'' for want of reliable swaption data. The mechanism belongs to the slow-moving-capital
literature: \citet{duffie2010} on the delayed arrival of arbitrage capital, \citet{bp2009} on the
margin--funding spiral, \citet{hkm2017} on intermediary capital as a pricing factor, and
\citet{sv1997} on the forced-unwinding channel that makes convergence trades risky. The
clientele-and-segmentation reading connects to \citet{vv2021} on preferred habitat and to
\citet{ks2023} on pension-driven demand in the swap market, which is the mechanism behind the sterling
result in Section~\ref{sec:within}. On the option side, \citet{cfj2021} document the variance risk
premium in Treasury futures, which is the object our third venue measures.

\section{Data}\label{sec:data}

\paragraph{One hub, read directly.} All Bloomberg series live in a single workbook,
\texttt{data/raw/Bloomberg/bbg\_paper2.xlsx}, with four sheets read directly by the loaders --- no
derived CSVs and no migration step, so there is exactly one place where a datum can be wrong.
\emph{\texttt{swap}}: 108 mid series, six curve families $\times$ the tenor grid $1$--$50$Y, now
including the $4$Y, $6$Y, $8$Y and $9$Y nodes added for the forward-rate factor.
\emph{\texttt{govt}}: the German, U.K.\ and Japanese government curves (\texttt{GDBR}, \texttt{GUKG},
\texttt{GJGB}) on the tenor grid $1$--$30$Y from 1995--1999; the long gilt and JGB futures
(\texttt{G 1}, \texttt{G 2}, \texttt{JB1}, \texttt{JB2}); the long-maturity government total-return
indices used for the term factors (\texttt{LUTLTRUU}, \texttt{I01656EU}, \texttt{I02667GB},
\texttt{I02988JP}, all unhedged in local currency, plus the aggregates \texttt{I00715JP} and
\texttt{I03374GB}); and the one-month sovereign bills (\texttt{GB1M}, \texttt{GETB1},
\texttt{UKGTB1M}). \emph{\texttt{cds}}: 22 dealer and index CDS series from 2001.
\emph{\texttt{rates \& vol}}: MOVE, VIX, V2X and the funding legs (\texttt{US0003M},
\texttt{USSOC}, \texttt{EUR003M}, \texttt{XESTR3M}, \texttt{EESWEC}, \texttt{USOSFRC}) from 1995.

\paragraph{Swaption surfaces.} Barclays Live, four batch files, 176 raw series filtered to 166 on
staleness and coverage; expiries $1$M--$1$Y $\times$ tails $2$Y--$30$Y, normal and lognormal families.

\paragraph{Government term structure.} Gürkaynak--Sack--Wright fitted zeros (\texttt{feds200628}),
\texttt{SVENY01}--\texttt{SVENY30} plus the native forward series, from 1985 --- true zero-coupon
yields rather than the par approximations used for swaps.

\paragraph{Term premia.} \texttt{ACMTermPremium.xls} (New York Fed, monthly and daily,
\texttt{ACMTP10}, from 1961) as the primary control, with Kim--Wright (\texttt{feds200533},
\texttt{THREEFYTP1000.B}, from 1990) as robustness. Both are U.S.-only; for the other markets a
Cochrane--Piazzesi-style forward factor is built from the local tenor grid
(Section~\ref{sec:termpremium}).

\paragraph{Factor libraries (from Paper~1).} \citet{fh2001} trend-following factors
(\texttt{TF-Fac.xls}, five PTFS lookback straddles, to March 2026); the \citet{hkm2017} intermediary
capital ratio; the NYU Stern illiquidity composite (daily, to June 2026); the
Jurado--Ludvigson--Ng macro, real and financial uncertainty series; the New York Fed primary-dealer
position statistics; and the OFR's DV01-weighted rendering of the CFTC Traders in Financial Futures
report (154 series, 2013--2026), which separates asset-manager, dealer and leveraged-fund positioning
in Treasuries and is the direct observation of the two clienteles the story posits.

\paragraph{Dealer composite.} U.S.\ leg: Morgan Stanley, JPMorgan Chase, Bank of America (Goldman
excluded: 2012 start and $29\%$ stale days). European leg: BNP, Deutsche Bank on the SLA clause with
full 2001 history, Barclays, UBS (the Deutsche SR series excluded: 2019 start). All constituents run
from 2001 with $3$--$8\%$ stale days.

\paragraph{Transaction costs.} Dealer-quoted full bid--ask on the relevant swap legs: $0.10$--$0.15\bp$
in calm conditions, $0.30\bp$ in medium, $1.00\bp$ in stress. This maps to a base half-spread of
$0.0625\bp$ on the $2$Y and $10$Y legs and $0.094\bp$ on the $30$Y, with a MOVE-tiered multiplier
$(1.0,\,2.4,\,8.0)$ at breakpoints $70$ and $110$. A \texttt{COST\_SAFETY} switch re-runs everything at
any multiple; results at $\times 2$ are reported.

\section{Methods}\label{sec:methods}

\paragraph{The $\sBE$ engine.} A uniform three-point construction on $2$/$10$/$30$Y, value- and
duration-matched, with $C=160$, applied identically to all eight markets --- four swap curves and four
government curves. The euro fly was moved from $2$/$10$/$25$ to $2$/$10$/$30$ once \texttt{EUSA30}
became available, both for uniformity and to remove an arbitrary-looking choice.

\paragraph{The C3 test.} Correlation of monthly changes and of levels between $\sBE$ and the swaption
ATM normal volatility, computed cell by cell across the cube and reported for the $3$M$\times$$10$Y
reference cell in the main text. Government curves are mapped to the swaption surface of their own
currency, exactly as the Treasury curve already was.

\paragraph{Costs and the two bounds.} Net returns are reported under two conventions: \emph{netA}
(patient, flip-only --- costs charged when the position changes sign) and \emph{netB} (full
replacement --- every leg re-traded each month), the latter being a deliberately pessimistic bound
that no practitioner would incur.

\paragraph{Risk-free convention.} Excess returns on the term factors subtract the local one-year OIS
rate, applied uniformly to all four currencies. This departs from Paper~1, which uses the one-month
sovereign bill, and the departure is forced: the U.K.\ bill series begins only in October 2008 and no
usable Japanese bill exists, so the Paper~1 convention would make the four term factors
non-comparable --- and cross-market comparability is the paper. The OIS is also the natural risk-free
in a rates setting, being the collateral and discounting rate for the instruments traded. A
\texttt{RF\_MODE} switch reproduces the bill convention where the data allow.

\paragraph{Currency convention.} The strategy return is a local-currency basis-point P\&L on a
value-matched, duration-neutral position: it is self-financing and carries no currency exposure. Term
factors are therefore built from indices \emph{unhedged in local currency} and no Glück--Hübel--Scholz
translation is applied --- a deliberate departure from Paper~1, where an EUR-based investor requires it.

\section{Results}\label{sec:results}

\subsection{The object across eight markets}\label{sec:object}
Mean $\sBE$ over each market's full sample: USD swap $89$, EUR $107$, GBP $79$, JPY $32$, Bund $109$,
gilt $66$, JGB $62$, UST $119\bp$/yr. Topological validation against the episode signs expected from
theory passes in $11$ of $12$ cells; the single exception is the euro in October 2022, where the
curve reads $+90$ (rich) against an expected neutral --- consistent with the euro area not being where
the 2022 balance-sheet constraint bound.

\subsection{The four questions of the antecedent paper}\label{sec:fourq}
\citet{rr2021} close with four questions on the \emph{relative cost} of trading convexity: what it
would cost to reproduce the long-gamma profile with a straddle of long-dated swaptions; whether the
implicit premium from the negative carry exceeds the option premium; which venue is more efficient at
a given moment; and whether a delta-, gamma- and vega-neutral package would carry a non-zero cost.
They state they lack reliable swaption prices and leave the questions open.

Since $\tilde\sigma_{BE}$ and $\sigma_{IV}$ are both normal volatilities of the same rate in basis
points per annum, the wedge $W_t=\tilde\sigma_{BE,t}-\sigma_{IV,t}$ is the relative price of
volatility across venues, and a vega-neutral package long one venue and short the other carries a
cost proportional to $W$ by construction. At equal volatility exposure the cost ratio is
$(\tilde\sigma_{BE}/\sigma_{IV})^2$, independent of the fly's convexity constant: $1.04$ (USD swap),
$1.81$ (EUR), $0.94$ (GBP), $0.41$ (JPY). The wedge reverses across regimes --- in the dollar it runs
$+106$\,bp through the crisis and $-131$ through the inflation episode; in sterling $+109$ to $-157$
--- and $|W|$ averages $83$--$106\%$ of the option premium with $t$ between $6.5$ and $9.8$ and
first-order autocorrelation of $0.95$--$0.99$.

\emph{What these numbers do and do not establish.} They answer the \emph{practical} question: a desk
pays quoted prices, premium included, so the raw ratio is the correct answer to ``where is convexity
cheaper.'' They do \emph{not} establish segmentation, for the reason set out in
Section~\ref{sec:vrp}. The explicit assembly of the neutral portfolio requires money premia and the
greeks and is outside the scope of a volatility-based measurement.

\subsection{The volatility risk premium, and the verdict on levels}\label{sec:vrp}
$\sigma_{IV}$ is a risk-neutral quantity; $\tilde\sigma_{BE}$ is the break-even for a realised-volatility
strategy and therefore physical. A wedge is expected even under perfect integration, so any
level-based claim is confounded. Decomposing
\[
W \;=\; \underbrace{\bigl(\tilde\sigma_{BE}-\mathbb{E}[RV]\bigr)}_{\text{residual }R}\;-\;\underbrace{VRP}_{\sigma_{IV}-\mathbb{E}[RV]},
\]
with realised volatility over the option horizon as the proxy for $\mathbb{E}[RV]$, gives a VRP of
$-1.3$ (USD, $t=-0.3$), $12.6$ (EUR, $t=4.0$), $5.9$ (GBP, $t=1.6$) and $19.0$\,bp (JPY, $t=6.1$).
The residual is $4.9$ ($t=0.3$), $37.8$ ($t=3.1$), $3.5$ ($t=0.2$) and $0.6$ ($t=0.1$). In three of
four markets the wedge is \emph{entirely} the premium: the yen's cost ratio of $0.41$ becomes $1.04$
once purged.

The level was then tested from seven further directions, and the verdict is uniform.
\emph{Construction of the realised leg}: windows of 21, 63 and 126 days, weekly sampling, and the
forward $3$M$\times$$10$Y rate --- the swaption's correct underlying --- move the euro residual only
between $36.8$ and $39.1$ and leave the others at zero, so the constant-maturity convention is not
the issue. \emph{Ex-ante forecasting}: an out-of-sample HAR-RV and $\sigma_{IV}$ purged of an
expanding-window VRP give the same answer as the ex-post construction on common dates (differences
of $1$--$2$\,bp), so forecast error is not what drives the result. \emph{Crises}: excluding all four
stress windows leaves the euro at $36.6$ ($t=2.8$) and the others at zero. \emph{Regimes}: in the
dollar and sterling the purged premium is $+97$ and $+89$\,bp through the zero-rate years and $-135$
and $-160$ through the inflation episode, with endogenous break dates of March 2022, March 2022 and
October 2021 --- so there is no ``level'' in those markets, only regimes that cancel.
\emph{Conditional sorting}: the purged premium sorts strongly on intermediary capital
unconditionally ($t=4.8$ and $4.4$), but the sign \emph{reverses} within sub-periods ($+70$ before
the break, $-101$ after, in the dollar), so the unconditional sort is largely ordering time.
\emph{Two-channel identification}: entering volatility stress and balance-sheet stress jointly gives
signs opposite to the mechanism's prediction, and coefficients estimated before the break predict
$+37$ and $+62$ where the realised post-break means are $-139$ and $-155$.

\emph{The conclusion, stated plainly.} The purged level is informative in the euro alone, where it
is robust across four forecasters, two samples and the exclusion of crises, and where it appears in
the Bund as well ($26.2$\,bp, $t=2.0$) --- a euro-area fact rather than a swap-market one, with the
paired swap-minus-government difference of $11.6$\,bp ($t=3.7$) as the instrument-specific layer. In
the other markets the level is not identified and is reported as such. This is a genuine narrowing
of what the paper claims relative to earlier drafts, and it is the reason the weight of the argument
rests on co-movement and on conditional returns, neither of which is a level statement.

\emph{One methodological caution for the write-up.} Because $R$ contains realised volatility with a
negative sign, any state variable correlated with realised volatility loads mechanically negative on
$R$. Regressions of the purged level on volatility-linked stress measures are therefore not
informative about the mechanism, and are not used as such.

\subsection{Segmentation: the central fact, and the backbone of the paper}\label{sec:c3}
Unlike the level, this is a statement about \emph{changes}, and a risk premium explains a level, not
the absence of correlation between monthly changes in two prices. That immunity was verified rather
than asserted: decomposing $\sigma_{IV}$ in real time into an expected-volatility component and a
premium, the changes in $\tilde\sigma_{BE}$ track \emph{neither} --- correlations with the expected
component are $-0.01$, $0.11$, $0.01$ and $0.08$, and the partial correlation with
$\Delta\sigma_{IV}$ controlling for the premium is identical to the raw one in every market. A
time-varying volatility risk premium therefore cannot manufacture the disconnect.
For the $3$M$\times$$10$Y reference cell, the correlation of monthly changes between $\sBE$ and $\sIV$
is $-0.01$ (USD swap), $+0.03$ (UST), $+0.24$ (EUR), $+0.18$ (Bund), $-0.07$ (GBP), $+0.07$ (gilt),
$+0.31$ (JPY), $+0.14$ (JGB). Two venues, one object, essentially disconnected.

Three robustness results matter more than the point estimates. First, the \textbf{full cube}: scanning
all $16$ cells ($4$ expiries $\times$ $4$ tails) in each of five markets, $56$ of $80$ ($70\%$) have
$|\rho|<0.20$; USD is $16/16$ with median $-0.00$, gilt $16/16$ with median $-0.07$, UST $16/16$ with
median $+0.06$. Critically, the ten-year tail column is \emph{not} systematically higher than the
others, which excludes the objection that we have compared mismatched objects. Second,
\textbf{stability}: on a rolling 60-month window the trend per decade is $-0.35$ (USD), $-0.43$ (GBP),
$-0.21$ (UST), $+0.15$ (EUR), $+0.08$ (JPY), with terminal values in line with the medians --- the
anomaly is not closing, and in the anglophone markets it is widening. Third, a \textbf{third venue}
(Section~\ref{sec:venues}).

\subsection{Which price carries volatility information}\label{sec:mz}
The segmentation result can be restated as a claim about \emph{information} rather than about
co-movement, and in that form it is positive rather than null. Regressing realised volatility over the
following quarter on each price, across all eight curves, the option's price carries slopes of
$0.73$ to $1.09$ with $t$ from $2.0$ to $5.0$, while the curve's price carries slopes near zero in six
of eight markets and significantly negative in sterling ($-0.11$, $t=-2.3$). Japan inverts on both of
its curves: the swap curve forecasts with slope $0.14$ ($t=4.6$, $R^2=0.30$) and the JGB curve with
$0.10$ ($t=4.8$), while the option term loses significance.

\emph{The claim must be relative, and here is why.} Out of sample against a trailing-realised-volatility
benchmark, neither price beats the naive forecast --- Campbell--Thompson $R^2$ is negative for the
option in every market --- and once trailing realised volatility enters the regression, the option's
coefficient becomes insignificant everywhere ($t$ between $0.1$ and $1.1$). Implied volatility is a
biased forecast and trailing realised volatility is a demanding benchmark; both facts are known. So
the paper does \emph{not} claim that the option forecasts volatility.

What it claims is the encompassing result, which is exactly the relative statement the argument needs.
Regressing the option's forecast error on the curve's price gives coefficients indistinguishable from
zero in every market except Japan, while regressing the curve's forecast error on the option's price
gives coefficients of $0.72$ to $1.09$ with $t$ up to $4.6$: \textbf{the option's price encompasses the
curve's, and the curve adds no residual information about future volatility}. This holds in five of
eight markets, with two indeterminate and Japan reversed. One of the two prices of volatility is
therefore not an information-bearing forecast relative to the other --- it is set by something else,
which is what the clientele reading asserts, and it gives quantitative content to the ``too much or too
little curvature'' of \citet{rr2021} and to the forecasting limitation \citet{rp2018} diagnose in their
own work.

\subsection{Curve shape, unspanned volatility, and what the object is}\label{sec:spanning}
The sharpest objection to the design is not the risk premium but the nature of $\tilde\sigma_{BE}$: it
is derived from the fly's carry and therefore from the curve's local slopes, so a referee may say the
object is curve shape rather than a volatility price, and that its failure to track implied volatility
is the known unspanned-volatility finding of \citet{cdg2002} and \citet{ab2010} rather than
segmentation. The objection deserves a direct answer, and the data give one.

Regressing changes in each price on changes in the level, slope and curvature of its own curve,
$\tilde\sigma_{BE}$ is spanned with $R^2$ of $0.50$ to $0.70$ --- substantially, as its construction
implies --- while $\sigma_{IV}$ is spanned with $R^2$ of only $0.15$ to $0.32$. Neither object is fully
spanned, which places both inside the unspanned-volatility literature rather than on opposite sides of
it.

The decisive test purges $\tilde\sigma_{BE}$ of the three curve factors and correlates the residual
with changes in the option's price. The correlations are $-0.06$, $+0.09$, $-0.06$ and $-0.04$: the
disconnect \emph{survives}, and in fact becomes uniform, since the euro and Japanese co-movement seen
in the raw series ($0.24$ and $0.31$) is attributable to the spanned component and vanishes once it is
removed. It is therefore not curve shape that fails to co-move with the option; it is the part of the
curve's volatility price that is \emph{not} curve shape.

This reframes the contribution rather than defending it. $\tilde\sigma_{BE}$ is, by construction, the
volatility that the curve's shape implies under no arbitrage --- the analogue of implied volatility, in
the antecedent paper's own words, not a second measurement of it. The finding is that this implied
quantity is set by flow rather than by expectation, and the paper's position with respect to the
unspanned-volatility literature is constructive: that literature establishes that volatility risk is
not spanned by the yield curve; this paper prices the wedge that non-spanning leaves open, and gives it
a limits-to-arbitrage foundation.

\subsection{Within-currency identification}\label{sec:within}
For each currency there are now two curves --- the swap curve and the government curve --- priced
against \emph{one} swaption surface. Institutional differences between countries are held fixed by
construction, so what varies is the instrument and its clientele.

\emph{The result, at the strength the data support.} In levels the contrast is large: sterling
$-0.61$ on the swap curve against $-0.44$ on the gilt curve, Japan $+0.65$ against $+0.73$. But
levels are the object Section~\ref{sec:vrp} declares confounded, and the honest test is in changes
with inference that accounts for the two curves' dependence (their changes correlate $0.55$--$0.81$).
Under a Williams paired test the swap-minus-government difference reaches significance only in
Japan ($+0.19$, $p=0.040$) --- and there with the \emph{opposite} sign to the level contrast ---
while sterling is marginal ($-0.11$, $p=0.066$) and the dollar and euro are indistinguishable from
zero.

What does survive is the \emph{systematics}: across all sixteen cells of the surface the sign of the
difference is uniform --- negative in all sixteen for sterling, positive in all sixteen for Japan,
with means of $-0.14$ and $+0.20$. A sign that is stable across the cube is not noise, but it is
weaker than a significant paired difference, and this section is therefore presented as corroborating
structure rather than as the paper's identification. The claim that it was the cleanest
identification does not survive its own test, and is withdrawn.

\subsection{Two exogenous clientele shifts, and why they do not identify}\label{sec:natexp}
Yield-curve control in Japan (September 2016 to March 2024) and the September--October 2022 gilt
episode are dated exogenously and shift the clientele by decree, which makes them the natural place
to look for identification. They do not deliver it, and the failure is reported because it disciplines
the mechanism section.

Around the YCC window the co-movement in levels is flat ($p=0.87$), and in changes the government
curve moves in the direction \emph{opposite} to the prediction ($-0.24$, $p=0.038$: co-movement falls
during the regime). The Japanese inversion in Section~\ref{sec:mz} rises monotonically across the
three sub-periods rather than peaking inside the regime, so it is not a YCC artefact --- and by the
same token yield-curve control cannot be offered as its cause. Around the gilt episode the predicted
attenuation is not significant in changes. Decisively, the placebo fails: the dollar and Treasury
curves move at the Japanese and British dates ($p=0.002$ to $p=0.039$), which is the signature of
synchronised global regime change rather than of a clientele shift. No causal claim is made from
either episode.

\subsection{The option side: a named clientele, and a term structure to test}\label{sec:optionside}
The economics of this paper names a clientele on the curve side --- liability-driven investors, mortgage
hedgers, official holders --- but leaves the option side generic, as desks whose vega capital binds.
That asymmetry is a weakness, and there is a specific candidate that closes it.

In the dollar market, servicers and holders of agency mortgage-backed securities are short the
prepayment option and therefore short convexity, and they hedge by buying receivers and Bermudan
swaptions. That demand is a duration-and-convexity clientele operating on the \emph{option} side, and it
is close to absent in Europe and Japan, where prepayable fixed-rate mortgage stock is small. This places
the swaption market's risk premium inside an existing literature --- \citet{duarte2008} on refinancing
activity and interest-rate volatility, \citet{perlisack2003} and \citet{hanson2014} on mortgage
convexity, and \citet{mmvv2016} on mortgage risk and the yield curve --- which the design has not so far
engaged.

The prediction is sharp and testable on data already in hand. If prepayment and Bermudan hedging shapes
the option premium, that premium should have a \emph{term structure across expiries and tenors} rather
than a single level, and the structure should be most pronounced in the dollar and least in Japan. A
fragment consistent with this is already in the robustness matrix: measured across the surface, the
wedge grows monotonically with option expiry in every market --- from $3$ to $11$ basis points in the
euro and from $-25$ to $-20$ in the yen on the baseline geometry --- which is a structure, not a level.
The full test is the premium itself, cell by cell across the sixteen expiry-tenor combinations, compared
across currencies with and without a prepayable mortgage market. It is the one place where the paper can
say something about \emph{why} the option premium is what it is, rather than only removing it.

The premium has now been measured cell by cell, and the surface is strongly structured --- but not in
the dimension the hedging hypothesis predicts.

\emph{The structure is in the tail, not the expiry.} Regressing each cell's mean premium on the logs of
expiry and tenor, the tenor slope dominates: $-16.6$ in the dollar, $-13.7$ in sterling, $-9.9$ in the
euro, with $R^2$ between $0.90$ and $0.99$. The dollar premium runs from $+24$bp at the two-year tail to
$-18$bp at the thirty-year tail, monotonically, with $t$ of $6.3$ and $-5.4$ at the two ends. Expiry
slopes are small and negative everywhere ($-1.7$ to $-4.1$), so the premium mildly \emph{declines} with
option expiry rather than rising.

\emph{Japan is the exception, and it is flat.} Its tenor slope is $-0.8$ with $R^2$ of $0.57$: the
premium is $15$--$19$bp at every tail from two to thirty years. Where the other three markets price
volatility risk very differently across the curve, Japan prices it almost uniformly --- which is what a
single dominant official participant compressing the premium would produce, and it is the one place where
the market with an atypical marginal investor also has an atypical premium surface.

\emph{The mortgage channel, as specified, is not supported, and this is reported rather than reframed.}
The prediction was a steeper expiry structure in the dollar. Pooling the sixty-four cells with a dollar
dummy interacted with log expiry, the interaction is $+1.6$ with $t=0.6$; sterling, not the dollar, has
the steepest expiry slope. What the data do show is a dollar--Japan contrast in the \emph{tenor}
dimension --- the dollar's tail slope is the steepest of the four and Japan's is one twentieth of it ---
but that is a different fact from the one the hypothesis named, and attributing it to prepayment hedging
would be fitting the mechanism to the result. The hypothesis is treated as any other: tested as stated,
and reported as not supported in that form.

\emph{Two methodological consequences, one of which reaches backwards.} First, the euro residual is
$38$ to $42$bp across all four expiries, so that result is robust to the horizon over which the realised
leg is measured as well as to the four constructions of Section~\ref{sec:vrp}. Second, and less
comfortably: the reference cell used throughout the earlier panels sits close to where the dollar premium
crosses zero. The dollar premium is $-1$bp at three months into ten years but $+24$bp at three months
into two years. The statement that the dollar's premium is slightly negative is therefore a property of
the chosen cell rather than of the market, and every level claim conditioned on a single cell inherits
that dependence. The comparison against the curve remains correct at the ten-year tail, because that is
the tail the fly's body matches; but the surface must be reported, not summarised by one of its points.

\emph{A caveat on power.} Long-tenor, long-expiry cells have few non-overlapping observations --- a
two-year cell has of the order of $T/24$ --- and the negative thirty-year premia may partly reflect the
2022 realisation. Those cells are reported with the standard errors that construction allows and are not
used to carry an argument.

\subsection{The mechanism}\label{sec:mechanism}
Sorting timed returns into terciles of dealer-CDS stress, the HIGH cell is significant in four of four
swap markets: USD $-10.30 \to 2.52 \to +19.64$ ($t=2.5$), EUR $4.90 \to 6.74 \to +21.23$ ($t=2.6$),
GBP $-20.66 \to 1.48 \to +23.44$ ($t=2.4$), JPY $2.34 \to 1.39 \to +9.66$ ($t=2.0$). Excluding all four
crisis windows the HIGH cell remains significant in three of four (USD $[2.0]$, EUR $[2.3]$,
JPY $[2.2]$; GBP falls to $[1.3]$ and is flagged): the mechanism is not an artefact of the episodes.
Using the senior-financials-minus-main differential instead of the composite, EUR is $[2.8]$ and GBP
$[2.5]$.

The third lens is decisive for the objection that the stress measure is our own construction. Sorting
on the \citet{hkm2017} intermediary capital ratio --- the canonical academic measure, independent of
our CDS data --- the premium concentrates where intermediaries are \emph{constrained} in four of four
markets: USD $+23.9$ ($t=2.1$), EUR $+19.0$ ($t=2.5$), GBP $+38.7$ ($t=3.4$), JPY $+8.2$ ($t=1.7$),
with low-minus-high spreads up to $+44\bp$.

Reversal half-lives (S5) order as the story predicts, with the most integrated market reverting
fastest: JPY $12.5$, USD $18.8$, EUR $20.3$, GBP $24.4$, UST $28.1$ months.

\subsection{Three conditional corollaries}\label{sec:corollaries}
A volatility-managed version of the strategy (Moreira--Muir) \emph{destroys} the premium in the
dollar, euro and sterling --- the euro's Sharpe falls from $0.64$ to $0.09$, with a managed alpha of
$-6.9$ ($t=-2.9$) --- because returns concentrate precisely in high-variance months. That is what
stress-state compensation implies, and it is the opposite of the equity-market pattern; it also
means the strategy cannot be volatility-targeted, which is stated rather than hidden. Japan, once
again, inverts ($+5.7$, $t=2.2$). A Ferson--Schadt specification with predetermined stress gives
directionally consistent but statistically weak loadings (three of four positive, sterling
significant at $t=2.0$), matching the real-time analogue's honest reading below: supporting
evidence, not a headline. And the wedge's persistence rises toward a unit root in high-stress months
in the dollar and sterling --- reversion effectively stops when balance sheets are constrained ---
while the euro shows no difference and Japan the opposite.

\subsection{The real-time analogue}\label{sec:conditional}
The conditional ladder uses in-sample thresholds and is therefore evidence about the mechanism, not an
implementable strategy. Its real-time analogue --- a rolling five-year threshold, no look-ahead ---
delivers conditional returns two to four times the unconditional in four of four markets (USD $+14.2$,
EUR $+19.0$, GBP $+9.5$, JPY $+10.3$), but the shift $p$-values are $0.083$--$0.269$. This is reported
as a supporting exhibit, not a headline.

\subsection{The third venue}\label{sec:venues}
The thirty-day at-the-money implied volatility on the ten-year bond futures (RX1 for the Bund, TY1 for
the Treasury, from 2004) prices the same volatility in a market distinct from both the curve and the
swaption surface. The test is whether the two \emph{option} venues agree with each other while both
disagree with the curve --- a pattern a noisy $\sBE$ could not produce. In USD the two option venues
correlate $+0.30$ with each other while the curve sits at $-0.01$ (swaptions) and $+0.09$ (futures); in
EUR the option venues correlate $+0.44$ and the curve $+0.24$ with both. The ordering is as predicted
in both currencies, sharply in USD. It is reported as corroboration and not as proof: the
option-to-option correlation is itself moderate, because a $3$M$\times$$10$Y swaption on the swap rate
and a one-month option on the bond future are genuinely different instruments. Extending the test to
gilts and JGBs requires their futures implied volatilities, which are not yet in the data.

\subsection{Not a known factor}\label{sec:factors}
The first paper's universe of tradable candidates supplies the sharpest available controls, and the
convexity strategy passes all of them.

\emph{Term.} Against the native term factors --- long-maturity government total-return indices,
unhedged in local currency, less the local risk-free --- alpha is $+7.40$ (USD, $[1.7]$), $+10.73$
(EUR, $[2.7]$), $+8.65$ (GBP, $[1.7]$), $+4.75$ (JPY, $[1.6]$) with $R^2$ of $0.00$--$0.02$. The
strategy is not a term premium in disguise.

\emph{Curvature.} Against \texttt{CURV\_2S5S10S}, the DV01-neutral futures butterfly --- the closest
tradable relative of the trade itself --- alpha is $+7.52$, $+9.89$, $+8.90$, $+4.97$ with $R^2$ of
$0.00$--$0.02$. If the strategy were curvature timing, this regressor would have absorbed it.

\emph{The whole rates block.} With term, curvature, slope and level entered jointly, alpha is $+7.59$
(USD, $[1.5]$), $+10.08$ (EUR, $[2.3]$), $+8.47$ (GBP, $[1.6]$), $+7.34$ (JPY, $[2.4]$), $R^2$
$0.07$--$0.15$.

\emph{Trend-following.} Against the five \citet{fh2001} PTFS lookback straddles --- including
\texttt{PTFSBD}, the bond straddle, which is by construction the closest thing to buying rate
volatility --- $R^2$ is $0.00$--$0.06$ and alpha is essentially unchanged. In the euro the loading on
\texttt{PTFSBD} is significantly \emph{negative} ($-0.71$, $t=-2.7$): if anything the strategy is short
bond-straddle exposure, the opposite of what the objection presumes.

A collateral result deserves its own line. Regressing $\Delta\sBE$ on the tradable curvature factor
gives $R^2$ of $0.00$--$0.03$ in all eight markets: $\sBE$ is \emph{not} the butterfly's realised
return. It measures the ex-ante \emph{price} of convexity, which is precisely the distinction the paper
claims.

\subsection{Attenuation bias, closed with arithmetic}\label{sec:atten}
If $\tilde\sigma_{BE}$ were noisy, the near-zero correlation would be mechanical. Building the object
on four node triples applied uniformly gives a reliability ratio --- the correlation of changes
across constructions --- of $0.79$, $0.79$, $0.83$ and $0.66$. Spearman's correction therefore moves
the correlations essentially not at all ($-0.01$ to $-0.01$ in the dollar, $-0.07$ to $-0.08$ in
sterling). Instrumenting one construction with another gives a beta of $0.017$ ($t=1.0$) with a
first-stage $F$ of $351$. The decisive framing is the threshold: for attenuation alone to produce a
true correlation of $0.50$ from the observed dollar figure, reliability would have to be $0.0002$,
and $0.022$ in sterling --- two to three orders of magnitude below what is measured. The objection
does not need to be argued away; it fails arithmetically.

\emph{Note for the write-up}: the term-premium loading of Section~\ref{sec:termpremium} must
\emph{not} be used as the anti-attenuation defence. It establishes that $\tilde\sigma_{BE}$ contains
\emph{a} systematic signal, not that it contains a \emph{volatility} signal, and invoking it invites
the worse objection that the measure is a term premium in disguise.

\subsection{The term-premium vulnerability, and the answer}\label{sec:termpremium}
The single most serious threat to the paper was found in-house before a referee could raise it.
Regressing $\Delta\sBE$ on stress and on the change in the ACM ten-year term premium, the term premium
enters with $t$ of $5.6$--$7.7$ in all five markets, and the stress loading flips sign or dies. Building
a Cochrane--Piazzesi forward factor \emph{locally} from each market's own $1$--$10$Y grid --- so the
control is no longer a U.S.\ proxy --- reproduces the result: $t$ of $3.7$--$7.1$ in all four markets.
The finding is real and must be disclosed.

The answer has two parts. \emph{Methodological}: $\sBE$ and the forward factor are both curve-shape
objects, so a level regression that includes one as a control for the other is over-controlled --- it
removes exactly the variation to be measured, and a strong association between them is economically
expected, since both are compensation for bearing rate risk. \emph{Empirical}, and decisive: the clean
test is a double sort, which neutralises the term premium by construction rather than by
over-controlling. Within each tercile of the local forward factor, high-stress states still pay more
than low-stress states in four of four markets: USD $+19.6$ $[1.5]$, EUR $+20.3$ $[2.2]$, GBP $+33.8$
$[2.6]$, JPY $+12.3$ $[1.7]$. Holding the term premium fixed, the returns remain concentrated where
intermediary balance sheets are constrained. The defensible position is therefore that $\sBE$ shares
much of its \emph{level} variation with the term premium --- expected --- while the strategy's
\emph{returns} carry no term-premium exposure at all, which the factor regressions of
Section~\ref{sec:factors} confirm directly with $R^2\approx 0$.

\subsection{Balance sheet, not macro}\label{sec:p5}
Replicating the design of Table~XIX of Paper~1 on the convexity wedge: the dependent variable is the
first principal component of the wedges across seven markets ($71\%$ of variance), regressed on
intermediary proxies and, separately, on a macroeconomic placebo.

\begin{center}
\begin{tabular}{lcc@{\hskip 2em}lcc}
\toprule
\multicolumn{3}{c}{Panel A: intermediary} & \multicolumn{3}{c}{Panel C: macro placebo}\\
\cmidrule(r){1-3}\cmidrule(l){4-6}
variable & coef & $t$ & variable & coef & $t$\\
\midrule
HKM        & $-0.864$ & $-2.33$ & $\Delta$UM & $+0.466$ & $0.43$\\
LIBOR--OIS & $+1.094$ & $+3.32$ & $\Delta$UR & $-0.217$ & $-0.27$\\
Dealer CDS & $+0.343$ & $+0.89$ & $\Delta$UF & $+0.710$ & $+1.61$\\
ILLIQ      & $-3.982$ & $-6.45$ &            &          &      \\
\midrule
\multicolumn{3}{c}{$\bar R^2 = 0.584$} & \multicolumn{3}{c}{$\bar R^2 = 0.130$}\\
\bottomrule
\end{tabular}
\end{center}

The contrast is $4.5\times$, against $2.2\times$ in Paper~1 ($0.233$ vs.\ $0.106$). This also upgrades the mechanism of \citet{rr2021}, who explain their abnormal returns by correlating them with measures of market liquidity: a correlational narrative becomes a test with a placebo that must fail. The test is not
that Panel~A works but that Panel~A works \emph{and Panel~C does not}: a wedge driven by macroeconomic
fundamentals would have reversed the ranking.

The negative loading on illiquidity is not an anomaly but the two-channel structure of
Section~\ref{sec:theory} appearing in a single regression. The NYU composite is genuinely a measure of
illiquidity (2008 mean $171$ against $26$--$45$ in calm years), and it loads negatively: more
illiquidity, \emph{richer} convexity. Volatility and illiquidity stress spike together and belong to
the flight channel; funding and dealer-credit stress belong to the balance-sheet channel; they push the
wedge in opposite directions. Paper~1 could not observe this because its three bases live on one
channel only. \emph{Caveat}: the placebo is currently partial --- the long-horizon inflation
expectation and economic-policy-uncertainty series of the original design are not yet in the data.

\subsection{The structural model: rejected}\label{sec:structural}
The model of Section~\ref{sec:structural-theory} is estimated as follows. For each market the P\&L
process is characterised by its volatility and first-order persistence; the arbitrageur holds for
twelve months, pays a balance-sheet charge, and is forced out if cumulative losses breach a barrier;
the required premium $\mu^\ast$ solves $CE(\mu^\ast)=0$ under CARA utility, with one absolute
risk-aversion parameter common to all markets.

The model fails. Predicted against observed monthly premia: EUR $10.41 \to 4.2$, GBP $8.88 \to 6.3$,
USD $7.38 \to 5.4$, JPY $4.82 \to 2.5$. The rank correlation between predicted and observed is $+0.20$
to $+0.40$ ($p \geq 0.60$) and the cross-sectional $R^2$ is negative. The rejection is robust to the
barrier specification: a fixed barrier and a VaR-scaled barrier at $k=0.5,\,1.0,\,1.5$ times
$\sigma\sqrt{T}$ all give the same verdict. One of the three comparative statics has the wrong sign:
greater loss persistence \emph{lowers} the required premium instead of raising it.

The diagnosis is the result. Across markets, the premium tracks segmentation --- rank correlation
$+0.80$ with the negative of the C3 level --- twice as strongly as it tracks the risk the arbitrageur
bears (rank $+0.40$ with P\&L volatility). The decisive case is the euro: it has the \emph{lowest}
P\&L volatility of the three large markets ($62$ against $72$ and $76$) and the \emph{highest} premium
($10.41$). That is incompatible with compensation for mark-to-market risk and consistent with
segmentation.

Why the rejection is informative rather than a failure of the project: a model calibrated on a single
market cannot be rejected at all, because the risk-aversion parameter adapts to whatever level is
observed. Here the model is over-identified and can fail; failing, it tells us that the convexity
premium is not compensation for the arbitrageur's own mark-to-market risk, and that a correct
structural model must carry clientele and the balance-sheet constraint as \emph{parameters}, not
merely the price dynamics. A single-market setting could not have discovered that, which is the
methodological value of building the panel first.

\section{The Price of the Friction: A Cross-Venue Strategy}\label{sec:strategy}
The rejected structural model asks what premium a constrained arbitrageur would require; this
section asks the dual and more direct question. If the two venues price the same volatility
differently, an arbitrageur can sell convexity where it is dear and buy it where it is cheap. Does
that trade pay? The answer completes the segmentation argument from the other side: the wedge
persists not because no one perceives it but because harvesting it does not clear the cost of
implementation, and this section measures that cost. What was, in the prior round, the single
undischarged item --- the assembly of the vega-neutral portfolio requires money premia --- is
resolved here, because the premia are computed from the surfaces themselves.

\subsection{The trade and its accounting identity}\label{sec:strat-identity}
In the euro, where the curve prices $97$ against the option's $72$ in normal-vol terms, the trade is
to short the fly (receive the carry, pay realised convexity) and hold a gamma-matched
$3\mathrm{M}\times10\mathrm{Y}$ ATM straddle, delta-hedged (pay the premium, receive realised). When
the two legs are matched in gamma, realised volatility cancels by construction, and the profit-and-loss
is the variance wedge:
\[
R_{LS} \;=\; k\big(\sigma_{BE}^2-\sigma_{IV}^2\big)\;+\;2k\,\Delta\sigma_{IV}^2,
\qquad k=\tfrac12\,C\,\Delta t\times10^4,
\]
where $C$ is the fly's convexity constant and $\Delta t$ the horizon. The constant $k$ is not
calibrated: it is the engine's own constant, the same one that \emph{defines} $\sigma_{BE}$ through
$\sigma_{BE}^2=-2\theta/(C\,\Delta t)$. The drift is the wedge in variance; the residual risk is the
mark-to-market of the remaining vega. This is a basis trade, not a volatility position: nothing in
$R_{LS}$ depends on the level of realised volatility. A strangle would be the natural
away-from-the-money analogue, but the surfaces in hand are ATM only, so the strangle is not priceable
until the smile is loaded --- the same measurement gap that Section~\ref{sec:open} lists. The ATM
straddle is the gamma-pure instrument and the cleanest replication of the comparison.

\subsection{Four architectures, one verdict}\label{sec:strat-four}
The trade is built four ways, each a stricter test than the last, and the results are produced by
\texttt{36}, \texttt{36b}, \texttt{36c}, \texttt{36d} and \texttt{36f}.

\emph{Always-invested, ideal hedge (\texttt{36}).} With the delta-hedge assumed continuous and free,
the fixed-direction trade returns $7$--$12.5$ bp/month with $t$ up to $2.7$ gross. But the
decomposition is revealing: in the euro the return is almost entirely drift (the wedge, $5.3$ of
$7.1$ bp), while in the dollar and sterling most of it --- $7$ to $9$ of $12$ bp --- is the duration
residual of the fly leg rather than the wedge. Only the euro harvests segmentation cleanly.

\emph{Daily-hedged simulation (\texttt{36b}).} Replacing the accounting identity with a genuine daily
Bachelier simulation --- monthly-rolled straddle, strike at the entry forward, gamma-matched sizing,
daily delta rebalanced at the dealer half-spread --- the simulated trade tracks the ideal
(correlation $1.00$, means $1$--$2$ bp lower), confirming the identity was a good approximation. The
killer is the cost of the hedge itself: $4.5$--$6$ bp/month. Net of hedge and package, the trade is
zero to slightly negative.

\emph{Conditional, band-hedged (\texttt{36c}).} With hedging at the desk-realistic tolerance band and
an economic entry threshold --- enter only when the expanding-mean wedge exceeds the current
round-trip cost --- the all-in nets are USD $+0.04$, GBP $+0.56$, EUR $-4.12$, JPY $-4.53$ bp/month.
The longer expiry ($1\mathrm{Y}\times10\mathrm{Y}$) makes it worse: the residual vega costs more than
the saved gamma. The internal placebo is exact: Japan, where the purged wedge is null, is inside the
threshold in zero months.

\emph{Per-trade, Paper-1 architecture (\texttt{36d}, \texttt{36f}).} Carrying the signal-weighted
sizing, convergence exit and hysteresis of the antecedent paper's CDS-bond strategy onto the wedge,
the level-threshold variant ($|W|>10$ bp, entry into every market beyond threshold, exit at zero or
maturity, equal- and signal-weighted) gives a signal-weighted portfolio of $4.92$ bp/month, $t=2.2$,
Sharpe $0.42$, with SW beating EW as in Paper~1. The result is real but modest: $92$--$94\%$ of exits
are at maturity rather than at convergence, so the object is a chained carry rather than a
converging basis; $118$ of $285$ trades are on the short-option side, which is variance-premium
selling the paper has excluded from the front door; and against the antecedent paper's net Sharpe
ratios of $0.98$--$1.56$, the figure is half.

\subsection{Why the euro yields so few trades, and what that means}\label{sec:strat-persist}
The per-trade architecture on deviations from the wedge ($s=W-\theta$, $\theta$ the expanding mean)
produces few euro trades despite a large average wedge, and the reason is structural, not a coding
artefact. In the euro the wedge is a \emph{persistent level}: $\theta$ absorbs it, so $s$ oscillates
around zero and rarely crosses the threshold. This is the decisive disanalogy with the CDS-bond basis
of Paper~1, which \emph{converges} to zero and so offers a natural round-trip. The euro wedge does not
converge --- it sits above realised volatility across the sample, through regimes. With a persistent
object, a per-trade architecture admits only two regimes: a level-threshold that puts one permanent
position on (the always-in carry of \texttt{36}/\texttt{36c}, in-market $89$--$91\%$ of months) or a
deviation-threshold that trades thinly (\texttt{36d}). ``Many trades'' is impossible in either,
because the wedge does not cross zero many times; it rests above it. The level is harvested by staying
in, not by entering and exiting --- which is why the \emph{recurring} cost, not the entry cost, is the
variable that decides.

\subsection{Costs}\label{sec:strat-costs}
The rate legs use real dealer quotes: the package half-spread as in the antecedent cost layer,
MOVE-tiered from full bid-offer quotes, and the hedge crosses the same spread on a duration-weighted
basis, together $8$--$10$ bp/month. The option leg's crossing cost is the one input without a dealer
quote in hand. The only published figure for a swaption round-trip is $0.75\times$vega for the entire
open-and-close (arXiv:1810.02125; Duyvesteyn and de Zwart 2015 report gross returns on the same four
markets). The placeholder used, $0.5$--$1.0$ bp/yr of normal vol per side, is $1.3$ to $2.7$ times
that figure, i.e.\ conservative. The sensitivity is settled: the straddle roll costs $0.3$ to $2.3$
bp/month against $8$--$10$ for hedge and package, so the verdict does not depend on the swaption
spread --- it would take roughly $10$ bp/yr per side to overturn it, an order of magnitude above
anything documented.

\subsection{The verdict, and what it establishes}\label{sec:strat-verdict}
Four architectures --- always-in ideal, daily-hedged simulation, band-hedged conditional, and the
antecedent paper's per-trade design with signal-weighted sizing and convergence exit --- and none
clears the friction band with statistical significance. The wedge is worth $5$--$7$ bp/month; the
measured round-trip is $9$--$14$, on real dealer quotes for the rate legs and conservative bounds for
the option leg. The only $t$ above $1.8$ anywhere in the package is Japanese variance-premium selling,
which the paper excludes by its own front door and labels as such. This is not an absence of result:
it is the quantitative form of the limits-to-arbitrage thesis. Segmentation persists because
arbitraging it does not pay, and the paper can now say by how much. The no-trade band is the answer to
the demand for a profitable strategy stated in the negative --- there is no free lunch, and that is the
finding --- and it opens the demand for an economic account of who keeps the two prices apart, which
is the subject of the clientele map that the next round must build and test.

\section{Established versus Open}\label{sec:open}
\textbf{Established.} The two prices of the volatility of the long rate do not co-move at the
margin, uniformly across the swaption cube ($97$ of $128$ cells below $0.20$ in absolute value; under FDR, statistically zero in the dollar and sterling, small but nonzero in the euro and yen), and the disconnect is not attenuation
(Section~\ref{sec:atten}), not a time-varying risk premium (Section~\ref{sec:c3}), not a known
factor (Section~\ref{sec:factors}) and not closing (rolling trends of $-0.35$, $-0.43$, $-0.21$ per
decade in the anglophone markets). In three of four markets the curve's price carries no information
about future realised volatility while the option's does, and Japan inverts
(Section~\ref{sec:mz}). Segmentation differs within a currency across the two curves priced against
one surface (Section~\ref{sec:within}). The premium is concentrated in constrained states under
three independent lenses and survives the exclusion of every crisis window
(Section~\ref{sec:mechanism}), is orthogonal to the tradable term, curvature, slope and level
factors and to the Fung--Hsieh straddles (Section~\ref{sec:factors}), and its common component loads
on balance-sheet proxies but not on a macroeconomic placebo, by a factor of $4.5$
(Section~\ref{sec:p5}). The three actors the mechanism names are visible in positioning data with
the predicted signs (Section~\ref{sec:optionside}). In the euro, and only there, the curve prices
convexity above expected realised volatility robustly (Section~\ref{sec:vrp}). The curve-shape objection is answered directly: purged of level, slope and curvature, the curve's price still fails to track the option's, and the failure becomes uniform across markets (Section~\ref{sec:spanning}). The volatility risk premium is a structured surface rather than a level: its tenor slope is steep in the dollar, sterling and euro ($R^2$ of $0.90$--$0.99$) and flat in Japan (Section~\ref{sec:optionside}). The pooled premium survives inference that admits cross-sectional dependence ($t=3.1$ clustered by date, $2.9$ Driscoll--Kraay, permutation $p=0.000$) and extends to the government curves (eight curves, pooled $t=4.8$). Japan inverts on four
independent margins --- forecasting, the wedge's sign, volatility management, and stress-persistence
--- a coherent signature of the YCC-era clientele rather than four coincidences. The wedge is priced
inside a measured friction band: built as a cross-venue trade four ways, including the antecedent
paper's signal-weighted per-trade design, it returns $5$--$7$ bp/month gross against a round-trip cost
of $9$--$14$ bp/month, and no architecture clears the band with significance
(Section~\ref{sec:strategy}). The one profitable corner is Japanese variance-premium selling, which
the paper excludes by construction. The no-trade band is the quantitative form of the
limits-to-arbitrage thesis and discharges the vega-neutral-assembly item by computing the option
premia from the surfaces.

\textbf{Open, and declared as such.} The purged \emph{level} is not identified outside the euro; the
regime reversal in the anglophone markets is documented but not explained by the state variables
available, and the two-channel account fails out of sample. The structural model is rejected in its
pure-risk form and a version carrying the clientele as a parameter is not yet built. No causal identification is
available: the two exogenous clientele shifts do not deliver it and the placebo fails
(Section~\ref{sec:natexp}), so the mechanism rests on a cross-sectional ordering, conditional sorts and
positioning data rather than on an experiment. The within-currency contrast is corroborating structure,
not identification. The option side's premium is now measured across the surface and is strongly structured in the tail, but the prepayment-hedging hypothesis does not explain that structure in the form in which it was tested, and no alternative account of the tenor slope is yet established. Whether Japan's flat premium surface is attributable to official participation is asserted by coherence, not identified. The
macroeconomic placebo is partial: $5$Y$5$Y inflation and policy uncertainty are missing. The
reliability panel of Section~\ref{sec:atten} uses a simplified roll and must be regenerated through
the Section~\ref{sec:methods} engine per node triple before submission. The cross-venue strategy of
Section~\ref{sec:strategy} carries one placeholder: the swaption crossing cost is a conservative
literature bound rather than a dealer quote, though the sensitivity analysis shows the verdict is
insensitive to it. The economic map of which clientele dominates each market --- insurers under
Solvency~II on the curve side, mortgage-convexity hedgers on the option side, official participation
in Japan --- is the natural next step and is not yet built as a tested section: the direct
clientele proxies, when confronted with the data, move the wedge in only some markets and with signs
that do not uniformly match the initial hypothesis, so the account must be rebuilt around the
channels that survive rather than asserted. Only the S1--S6 battery is pre-specified: the cube
scan, the local forward factor, the attenuation panel, the volatility-premium decomposition and the
cross-venue strategy were run in response to objections and must be labelled as such rather than as
pre-registered.

\section{Register of Runs}\label{sec:register}
All results are produced by \texttt{convexity/}, run in order:
\texttt{check\_inputs} $\to$ \texttt{01\_import\_data} (loaders and the cost layer) $\to$
\texttt{02\_sigbe\_engine} (the object, eight markets) $\to$ \texttt{03\_c3\_segmentation} (the central
fact) $\to$ \texttt{04\_costs\_net} (net returns, two bounds) $\to$
\texttt{05\_conditional\_tercile} (the ladder and its ex-crisis panel) $\to$ \texttt{06\_make\_figures}
$\to$ \texttt{07\_story\_tests} (S2, S5) $\to$ \texttt{08\_robustness\_audit} (staleness, rank,
episode exclusion) $\to$ \texttt{09\_mechanism} (S1, the two channels) $\to$ \texttt{10\_surface} (the
128-cell cube scan) $\to$ \texttt{11\_conditional\_strategy} (real-time analogue, two threshold
specifications) $\to$ \texttt{12\_controls} (Fung--Hsieh alpha, HKM ladder, ACM term premium, double
sort) $\to$ \texttt{13\_tp\_stability} (local Cochrane--Piazzesi factor, rolling C3) $\to$
\texttt{14\_venues} (the third venue) $\to$ \texttt{15\_paper1\_factors} (term, curvature, slope,
level) $\to$ \texttt{16\_p5\_panel} (intermediary versus macro placebo) $\to$
\texttt{17\_structural} (the model and its verdict) $\to$ \texttt{18\_clientele} (the two clienteles read from CFTC positioning) $\to$ \texttt{19\_attenuation} (errors-in-variables across node triples) $\to$ \texttt{20\_rebonato\_questions} (the four questions of the antecedent paper) $\to$
\texttt{21\_vrp} and \texttt{22\_vrp\_battery} (the premium, measured and removed; eight-panel battery)
$\to$ \texttt{23\_wedge\_diag} (episode diagnostics) $\to$ \texttt{24\_exante} (ex-ante forecasters,
sample reconciliation) $\to$ \texttt{25\_regime} (era structure, endogenous break) $\to$
\texttt{26\_two\_channels} (joint identification, rejected) $\to$ \texttt{27\_robust\_matrix} (node
geometry by surface cell) $\to$ \texttt{28\_paper1\_tools} (Forbes--Rigobon, FDR, block bootstrap,
Moreira--Muir, best subset) $\to$ \texttt{29\_dynamic\_tools} (DCC-GARCH, quantile, Ferson--Schadt,
regime persistence) $\to$ \texttt{30\_spanning} (curve-factor spanning and the purged co-movement)
$\to$ \texttt{31\_c3\_delta} (within-currency in changes, Williams paired test) $\to$
\texttt{32\_forecast\_race} (Mincer--Zarnowitz, out-of-sample, encompassing) $\to$
\texttt{33\_natural\_experiments} (YCC and LDI, with placebo) $\to$ \texttt{34\_premium\_robustness}
(double sort, clustered and permutation inference, eight curves) $\to$ \texttt{35\_vrp\_surface} (the
premium cell by cell). Seed and parameters are fixed in
\texttt{config.py}; the cost multiplier is exposed as \texttt{COST\_SAFETY} and the risk-free
convention as \texttt{RF\_MODE}.

\section{Road to Completion and Publication Strategy}\label{sec:road}

\paragraph{What remains.} Every standard objection now has a designed and executed answer, and two
that were not standard --- the risk premium and curve-factor spanning --- have been answered as well.
What remains is of three kinds. One substantial piece of theory: a structural model carrying the
clientele and the balance-sheet constraint as parameters, the pure-risk form having been rejected.
Four measurement closures: the smile at away-from-the-money strikes, without which a segmentation
confined to those strikes stays invisible; the missing thirty-year tail; the gilt and JGB futures
implied volatilities, so the third venue covers four markets rather than two; and the two placebo
series. And two regenerations through the production engine rather than the simplified roll: the
reliability ratios and the node-geometry matrix, whose signs are established but whose levels are not
yet comparable with the production series.

\paragraph{What the last round of tests changed.} Three claims were withdrawn rather than defended:
the within-currency contrast as the paper's identification, any causal attribution of the Japanese
inversion to yield-curve control, and the absolute statement that the option forecasts realised
volatility. Two were strengthened: the central fact survives purging of the curve factors and becomes
uniform across markets, and the pooled premium survives inference that admits cross-sectional
dependence while extending to the government curves. The architecture of the argument follows from
this: co-movement and the encompassing result are load-bearing, the euro level is a section.

\paragraph{Target.} On the current package and without a working structural model, the \emph{Journal of
Financial Economics} and the \emph{Review of Financial Studies} are realistic: the fact is new, the
within-currency identification is strong, and the P5 contrast is sharper than in Paper~1. The
\emph{Journal of Finance} generally expects the empirics to be anchored in a model that works; with the
structural extension solved, the paper enters that conversation. There is also a third route worth
keeping open: the \emph{Journal of Financial Economics} publishes work that documents a fact and
\emph{rejects} the obvious explanation for it, when the rejection is as clean as this one --- in which
case the structural failure becomes a section rather than an obstacle.

\begin{thebibliography}{99}
\bibitem[Brunnermeier and Pedersen(2009)]{bp2009} Brunnermeier, M.\ K., and L.\ H.\ Pedersen, 2009,
Market liquidity and funding liquidity, \emph{Review of Financial Studies} 22, 2201--2238.
\bibitem[Cremers et al.(2021)]{cfj2021} Cremers, M., M.\ Halling, and D.\ Weinbaum, and related work on
Treasury variance risk premia, 2021.
\bibitem[Duffie(2010)]{duffie2010} Duffie, D., 2010, Presidential address: Asset price dynamics with
slow-moving capital, \emph{Journal of Finance} 65, 1237--1267.
\bibitem[Fung and Hsieh(2001)]{fh2001} Fung, W., and D.\ A.\ Hsieh, 2001, The risk in hedge fund
strategies: Theory and evidence from trend followers, \emph{Review of Financial Studies} 14, 313--341.
\bibitem[He et al.(2017)]{hkm2017} He, Z., B.\ Kelly, and A.\ Manela, 2017, Intermediary asset pricing:
New evidence from many asset classes, \emph{Journal of Financial Economics} 126, 1--35.
\bibitem[Klingler and Sundaresan(2023)]{ks2023} Klingler, S., and S.\ Sundaresan, on pension-driven
demand and the swap market.
\bibitem[Rebonato(unpublished)]{rebonatoNote} Rebonato, R., unpublished note: certainty-equivalent
entry condition and the mark-to-market barrier for a convergence trade.
\bibitem[Rebonato and Putyatin(2018)]{rp2018} Rebonato, R., and V.\ Putyatin, 2018, on the break-even
volatility construction.
\bibitem[Andersen and Benzoni(2010)]{ab2010} Andersen, T.\ G., and L.\ Benzoni, 2010, Do bonds span volatility risk in the U.S.\ Treasury market? A specification test for affine term structure models, \emph{Journal of Finance} 65, 603--653.
\bibitem[Collin-Dufresne and Goldstein(2002)]{cdg2002} Collin-Dufresne, P., and R.\ S.\ Goldstein, 2002, Do bonds span the fixed income markets? Theory and evidence for unspanned stochastic volatility, \emph{Journal of Finance} 57, 1685--1730.
\bibitem[Duarte(2008)]{duarte2008} Duarte, J., 2008, The causal effect of mortgage refinancing on interest-rate volatility: Empirical evidence and theoretical implications, \emph{Review of Financial Studies} 21, 1689--1731.
\bibitem[Hanson(2014)]{hanson2014} Hanson, S.\ G., 2014, Mortgage convexity, \emph{Journal of Financial Economics} 113, 270--299.
\bibitem[Malkhozov, Mueller, Vedolin, and Venter(2016)]{mmvv2016} Malkhozov, A., P.\ Mueller, A.\ Vedolin, and G.\ Venter, 2016, Mortgage risk and the yield curve, \emph{Review of Financial Studies} 29, 1220--1253.
\bibitem[Perli and Sack(2003)]{perlisack2003} Perli, R., and B.\ Sack, 2003, Does mortgage hedging amplify movements in long-term interest rates?, \emph{Journal of Fixed Income} 13, 7--17.
\bibitem[Rebonato and Ronzani(2021)]{rr2021} Rebonato, R., and R.\ Ronzani, 2021, Is convexity
efficiently priced? Evidence from international swap markets, \emph{Journal of Economics and Finance}
63, 392--413.
\bibitem[Shleifer and Vishny(1997)]{sv1997} Shleifer, A., and R.\ W.\ Vishny, 1997, The limits of
arbitrage, \emph{Journal of Finance} 52, 35--55.
\bibitem[Vayanos and Vila(2021)]{vv2021} Vayanos, D., and J.-L.\ Vila, 2021, A preferred-habitat model
of the term structure of interest rates, \emph{Econometrica} 89, 77--112.
\end{thebibliography}

\end{document}
